\documentclass[11pt]{article}
\usepackage[letterpaper,margin=1in]{geometry}
\usepackage[T1]{fontenc}
\usepackage{mathptmx}
\usepackage{microtype}
\usepackage{enumitem}
\setlength{\parindent}{0pt}
\setlength{\parskip}{6pt}
\pagestyle{empty}

\begin{document}

\begin{center}
{\large\bfseries Team Information and Relevant Background}\\[3pt]
{\normalsize Metal-Poor Fossils in the Galactic Midplane: A NEWFIRM Pathfinder for the Hidden Galaxy Explorer}
\end{center}

\textbf{Team expertise and roles.---}
David Nidever (PI; Montana State University) will lead the survey design,
observations, data reduction, catalog construction, and delivery of the target
list to the Hidden Galaxy Explorer (HGE). He is the HGE Program Head and has
extensive experience with Galactic stellar populations, APOGEE spectroscopy,
and large imaging surveys. He led the SMASH survey with Blanco/DECam,
co-developed the APOGEE data-reduction pipeline, and led construction of the
NOIRLab Source Catalog. This combination of wide-field survey, photometric
pipeline, and near-infrared spectroscopic experience directly supports the
proposed program.

Gail Zasowski (University of Utah), the HGE Survey Scientist, will lead the
calibration and assessment of the metal-poor-star selection. Her expertise
includes Galactic archaeology, infrared extinction, APOGEE target selection,
and the chemical structure of the Milky Way. She led the target-selection work
for APOGEE and APOGEE-2 and developed the Rayleigh-Jeans Color Excess approach
used to estimate extinction toward highly obscured Galactic fields. She will
connect the NEWFIRM photometry to the HGE targeting strategy and help quantify
the resulting selection function.

Aaron Meisner (NSF NOIRLab) will contribute expertise in near-infrared
photometry, large imaging data sets, and public survey products. His work on
the full-sky unWISE coadds and catalogs provides directly relevant experience
with background-limited infrared imaging, crowded-source photometry, survey
calibration, and reproducible data products. He will advise on the photometric
processing, catalog validation, and public release. Guy Stringfellow
(University of Colorado Boulder) will assist with the classical NEWFIRM
observations and real-time assessment of data quality. His background in
infrared stellar astrophysics and APOGEE spectroscopy provides additional
experience relevant to observing and validating Galactic stellar populations.

\textbf{Resources and execution.---}
The team has already assembled the HGE Input Catalog for the proposed
footprint, including deep near-infrared and mid-infrared photometry, extinction
estimates, and nearly one million candidate luminous giants. Synthetic
photometry for the NEWFIRM filters and existing APOGEE spectra provide the
foundation for calibration and validation. The PI has access to the computing
and storage needed for image processing, crowded-field photometry, catalog
construction, and cross-matching. Nidever and Stringfellow will support the
classical observing run; Nidever and Meisner will oversee image and catalog
processing; and Zasowski and Nidever will calibrate the metallicity selection
and produce the ranked HGE target list. The proprietary period will be waived,
and the calibrated photometric catalog and metal-poor candidate list will be
released publicly after validation.

\textbf{Previous NOIRLab allocations.---}
The PI has received no NOIRLab telescope-time allocations as PI during the
past two years.

\end{document}
